\begin{dang}{Bài toán nâng cao định luật Bôilơ - Mariốt}
	\setcounter{paragraph}{0}
	\Anlythuyet{
		\paragraph{Áp suất chất lỏng} 
		\begin{align*}
			\boder{p=p_a+dh=p_a+\rho gh}\quad \begin{cases}
				p_a:\ \ct{áp suất từ bên ngoài nén lên mặt chất lỏng}\ (Pa)\\
				p:\ \ct{áp suất ở đáy cột chất lỏng}\ (Pa)\\
				h:\ \ct{là độ sâu tính từ mặt thoáng chất lỏng}\\ 
				\quad \ct{đến điểm tính áp suất} \ (m)\\
				\rho:\ \ct{khối lượng riêng của chất lỏng}\ (kg/m^3)\\
				d:\ \ct{trọng lượng riêng của chất lỏng}\ (N/m^3)\\
			\end{cases}
		\end{align*}
		%\paragraph{Áp suất khí quyển}
		%\Textbox{0.65}{\begin{itemize}
				%\item Trái Đất được bao bọc bởi một lớp không khí dày tới hàng ngàn kilômét, gọi là khí quyển.
				%\item Con người và mọi sinh vật khác trên mặt đất đều đang sống “dưới đáy” của “đại dương không khí” khổng lồ này.
				%\item Do không khí cũng có trọng lượng nên Trái Đất và mọi vật trên Trái Đất đều chịu áp suất của lớp không khí bao bọc xung quanh Trái Đất. Áp suất này tác dụng theo mọi phương và được gọi là áp suất khí quyển.
				%\item Càng lên cao không khí càng loãng nên áp suất khí quyển càng giảm.
				%\end{itemize}
				%}
			%\Textbox{0.35}{\centering
				%\begin{tikzpicture}[scale=1,thick,>=stealth']
				%\tkzDefPoints{2/0/a, 2.55/0/b}
				%\draw[fill=blue!40,rounded corners=5](1,0)--++(-90:1)--++(0:3)--++(90:1);
				%\draw[rounded corners=5](1,0.2)--++(-90:1.2)--++(0:3)--++(90:1.2);
				%\draw[rounded corners=2,fill=white](0,4)--++(-90:5)--++(0:0.1)--++(90:5);
				%\draw[rounded corners=2,fill=red!50](0,4)--++(-90:5)--++(0:0.1)--++(90:5);
				%\draw[<->](0.5,4)--node[midway,fill=white]{$1m$}(0.5,-1);
				%\draw[rounded corners=2,fill=white](2.5,-0.2)--++(90:5)--++(0:0.1)--++(-90:5);
				%\draw[rounded corners=2,fill=red!50](2.5,-0.2)--++(90:4)--++(0:0.1)--++(-90:4);
				%\draw[<->](3.5,3.8)--node[midway,fill=white]{$76cm$}(3.5,0);
				%\draw[](0.2,4)--(0.8,4)(0.2,-1)--(0.8,-1)(3.2,3.8)--(3.8,3.8);
				%\draw(a)--++(90:0.5)--++(180:0.5)node[above right]{$A$};
				%\draw(b)--++(45:0.5)--++(0:0.5)node[above left]{$B$};
				%\draw(2.55,4.2)--++(45:0.5)--++(0:1)node[above left]{Chân}node[below left,yshift=2pt,xshift=2pt]{không};
				%\tkzDrawPoints[size=2,fill=white](a,b)
				%\end{tikzpicture}
				%}
			%\begin{itemize}
			%\item Để đo áp suất khí quyển người ta dùng ống Tô-ri-xe-li:
			%\begin{itemize}
			%\item Lấy một ống thuỷ tinh một đầu kín dài khoảng 1 m, đổ đầy thuỷ ngân vào.
			%\item Lấy ngón tay bịt miệng ống lại rồi quay ngược ống xuống.
			%\item Nhúng chìm miệng ống vào một chậu đựng thuỷ ngân rồi bỏ ngón tay bịt miệng ống ra, thuỷ ngân trong ống tụt xuống, còn lại khoảng h nào đó tính từ mặt thoáng của thuỷ ngân trong chậu.
			%\end{itemize}
			%\item Độ lớn của áp suất khí quyển bằng áp suất của cột thuỷ ngân trong ống Tô-ri-xe-li.
			%\item Đơn vị đo áp suất khí quyển thường dùng là mmHg.
			%\end{itemize}
			\begin{note}
				Vì áp suất khí quyển bằng áp suất gây ra bởi cột thủy ngân trong thí nghiệm Tô-ri-xe-li, nên người ta còn dùng chiều cao của cột thủy ngân này để diễn tả độ lớn của áp suất khí quyển.\\
				\textit{Ví dụ:} áp suất khí quyển ở bãi biển Sầm Sơn vào khoảng 76 cmHg (760 mmHg).\\
				Áp suất khí quyển trung bình ở mực nước biển bằng: $p_0=101300\ Pa$.\\	
				Dụng cụ để đo áp suất là áp kế.
		\end{note}	}
		
	\end{dang}